%% kapa: exam=1 round=33 year=2022 period=1 form=A subject=economics subject_ko=경제학원론 seq=5 number=45
%% source: Q-Net 감정평가사 기출문제 (cst00309) artlSeq=5209997 file=raw/1008/5209997/2022년도 제33회 감정평가사 1차시험 1교시(민법,경제학,부동산학).pdf
%% conversion: kordoc 4.24.51 -> tools/gen1.py (auto); review=proofread-sonnet; flags=none
\begin{question}{45}
( )에 들어갈 내용으로 옳지 않은 것은? (단, 수요곡선은 우하향, 공급곡선은 우상향한다.)
\begin{framedtext}
정부는 X재에 대해 종량세를 부과하려고 한다. 동일한 세율로 판매자에게 부과하는 경우와 구매자에게 부과하는 경우를 비교할 때, (\quad\quad\quad\quad)
\end{framedtext}
\begin{choices}
  \item 구매자가 내는 가격은 동일하다.
  \item 판매자가 받는 가격은 동일하다.
  \item 조세수입의 크기는 동일하다.
  \item 균형 거래량이 모두 증가한다.
  \item 총잉여는 모두 감소한다.
\end{choices}
\end{question}
\answer{4}
